% Blueprint chapter for VerifiedReserving/ODP.lean.
% Renshaw and Verrall, A stochastic model underlying the chain-ladder technique,
% British Actuarial Journal 4 (1998) 903-923.
% Mack, A simple parametric model for rating automobile insurance or estimating
% IBNR claims reserves, ASTIN Bulletin 21 (1991) 93-109.
% Same zero-based conventions as content.tex; the cell (i,k) is observed when
% i + k <= n-1, and d = n-1-i is the latest observed development year of row i.

\chapter{The over-dispersed Poisson model}

Renshaw and Verrall, \emph{A stochastic model underlying the chain-ladder technique},
British Actuarial Journal 4 (1998) 903--923, put the run-off triangle in incremental
form, $X_{i,k} = C_{i,k} - C_{i,k-1}$ with $X_{i,0} = C_{i,0}$, and model the
incrementals as an over-dispersed Poisson (ODP) generalized linear model: independent,
with mean $m_{i,k}$ and variance $\varphi\, m_{i,k}$, and a log link carrying a row
effect and a column effect, so the mean is multiplicative,
\[
  m_{i,k} = a_i b_k , \qquad \sum_{k=0}^{n-1} b_k = 1 .
\]
Mack, \emph{A simple parametric model for rating automobile insurance or estimating
IBNR claims reserves}, ASTIN Bulletin 21 (1991) 93--109, starts from the same
multiplicative mean and estimates it from the \emph{marginal sums}. Both papers arrive
at the same system: for each accident year $i$ and each development year $k$,
\[
  \sum_{k \,:\, i+k \le n-1} m_{i,k} = \sum_{k \,:\, i+k \le n-1} X_{i,k},
  \qquad
  \sum_{i \,:\, i+k \le n-1} m_{i,k} = \sum_{i \,:\, i+k \le n-1} X_{i,k}.
\]
These are the score equations of the quasi-Poisson log-likelihood
$\sum (X \log m - m)/\varphi$ in the row and column parameters; the dispersion
$\varphi$ multiplies the whole score and cancels, which is why the ODP fit is the
Poisson fit and why the estimates do not depend on $\varphi$.

The chapter proves the two halves of the classical equivalence. Chain ladder
\emph{solves} the system (Theorems~\ref{thm:odpRow} and~\ref{thm:odpCol}), and a
positive multiplicative solution of the system \emph{is} chain ladder
(Theorem~\ref{thm:odpUnique}).

\emph{Source note.} Neither paper was available in full text for this formalization.
What is formalized is the marginal-sum system as quoted above and as it is stated in
the secondary literature (England and Verrall, \emph{Stochastic claims reserving in
general insurance}, British Actuarial Journal 8 (2002) 443--518, Section 2.3;
W\"uthrich and Merz, \emph{Stochastic Claims Reserving Methods in Insurance} (2008),
Section 2.3), in the two papers' own terms. The displayed equations of Renshaw and
Verrall (Section 3) and of Mack (Section 2) are therefore identified by section and by
the display itself, not by an equation number. Nothing here is probabilistic: the ODP
assumption enters only through the shape of the score equations, which are algebraic
identities in the fitted values.

\section{Definitions}

\begin{definition}[Incremental claims]\label{def:incr}\lean{VerifiedReserving.incr}\leanok
$X_{i,k} = C_{i,k} - C_{i,k-1}$ for $k \ge 1$ and $X_{i,0} = C_{i,0}$.
\end{definition}

\begin{definition}[Observed row and column]\label{def:obs}\lean{VerifiedReserving.obsRow, VerifiedReserving.obsCol}\leanok
The cell $(i,k)$ is observed when $i + k \le n-1$. The observed development years of
accident year $i$ are $\mathrm{obsRow}(n,i) = \{0,\dots,n-1-i\}$ and the observed
accident years of development year $k$ are $\mathrm{obsCol}(n,k) = \{0,\dots,n-1-k\}$.
The second one at $k+1$ is the contributor set of $\hat f_k$
(\texttt{obsCol\_succ}).
\end{definition}

\begin{definition}[Chain-ladder fit on the observed triangle]\label{def:CLfit}\lean{VerifiedReserving.CLcum, VerifiedReserving.CLincr}\leanok
\uses{def:fhat, def:Chat, def:obs}
The fitted cumulative on an observed cell is the latest observed entry of the row
divided back by the development factors,
\[
  \hat C_{i,k} = \frac{C_{i,\,n-1-i}}{\prod_{j=k}^{n-2-i} \hat f_j},
\]
and the fitted incremental is $\hat X_{i,k} = \hat C_{i,k} - \hat C_{i,k-1}$ with
$\hat X_{i,0} = \hat C_{i,0}$. These are the $\hat C$ projections of
Definition~\ref{def:Chat} read backwards from the latest diagonal: equivalently
$\hat C_{i,k} = \hat C_{i,n-1} / \prod_{j=k}^{n-2}\hat f_j$, the chain-ladder ultimate
times the cumulative development pattern up to $k$
(\texttt{CLcum\_eq\_ultimate\_div}).
\end{definition}

\begin{definition}[Multiplicative model and its pattern]\label{def:mult}\lean{VerifiedReserving.multFit, VerifiedReserving.PatternNormalized, VerifiedReserving.patternCum}\leanok
$m_{i,k} = a_i b_k$; the identifiability constraint $\sum_{k<n} b_k = 1$; and the
cumulative pattern $B_k = \sum_{j \le k} b_j$.
\end{definition}

\begin{definition}[Score (marginal-sum) equations]\label{def:score}\lean{VerifiedReserving.RowTotals, VerifiedReserving.ColumnTotals, VerifiedReserving.ScoreEquations}\leanok
\uses{def:incr, def:obs}
The row marginal sums, the column marginal sums, and their conjunction: the fitted
values reproduce the observed row and column totals of incrementals on the observed
triangle.
\end{definition}

\begin{definition}[Quasi-Poisson log-likelihood]\label{def:qll}\lean{VerifiedReserving.rowQuasiLogLik, VerifiedReserving.colQuasiLogLik}\leanok
\uses{def:incr, def:obs}
$\sum_{k} \big(X_{i,k}\log(t\,b_k) - t\,b_k\big)$ as a function of the row parameter
$t$, and $\sum_{i} \big(X_{i,k}\log(a_i\,t) - a_i\,t\big)$ as a function of the column
parameter $t$ (dispersion $\varphi = 1$, which only scales the expression).
\end{definition}

\section{The score equations are score equations}

\begin{lemma}[Row and column scores]\label{lem:score}\lean{VerifiedReserving.hasDerivAt_rowQuasiLogLik, VerifiedReserving.hasDerivAt_colQuasiLogLik}\leanok
\uses{def:qll}
The derivative of the quasi-Poisson log-likelihood in a row parameter is
$\big(\sum_k X_{i,k} - \sum_k t\,b_k\big)/t$, and in a column parameter
$\big(\sum_i X_{i,k} - \sum_i a_i\,t\big)/t$.
\end{lemma}
\begin{proof}\leanok
Termwise: $\frac{d}{dt}\big(X\log(tb) - tb\big) = X/t - b$, then sum over the observed
cells.
\end{proof}

\begin{lemma}[Stationarity is the marginal-sum equation]\label{lem:scoreZero}\lean{VerifiedReserving.deriv_rowQuasiLogLik_eq_zero_iff, VerifiedReserving.deriv_colQuasiLogLik_eq_zero_iff}\leanok
\uses{lem:score, def:score, def:mult}
For nonzero parameters, the row score vanishes exactly at the row marginal-sum equation
and the column score exactly at the column one. This is what licenses the name "score
equations" for Definition~\ref{def:score}, and it is where the dispersion $\varphi$
drops out.
\end{lemma}
\begin{proof}\leanok
Lemma~\ref{lem:score} and $a/b = 0 \iff a = 0$ for $b \ne 0$.
\end{proof}

\section{Chain ladder solves the score equations}

\begin{lemma}[Telescoping]\label{lem:telescope}\lean{VerifiedReserving.sum_incr, VerifiedReserving.sum_incr_obsRow, VerifiedReserving.sum_CLincr}\leanok
\uses{def:incr, def:CLfit}
$\sum_{k \le m} X_{i,k} = C_{i,m}$ and $\sum_{k \le m} \hat X_{i,k} = \hat C_{i,m}$.
\end{lemma}
\begin{proof}\leanok Induction on $m$. \end{proof}

\begin{theorem}[Row totals]\label{thm:odpRow}\lean{VerifiedReserving.chainLadder_fitted_row_totals, VerifiedReserving.chainLadder_fitted_row_totals_eq}\leanok
\uses{lem:telescope, def:score}
For every accident year $i$,
$\sum_{k \le n-1-i} \hat X_{i,k} = C_{i,\,n-1-i} = \sum_{k \le n-1-i} X_{i,k}$: the
chain-ladder fit reproduces every observed row total. This is the first marginal-sum
family of Mack (1991), Section 2, and the row score equation of Renshaw and Verrall
(1998), Section 3. No hypothesis is needed.
\end{theorem}
\begin{proof}\leanok
Both sides telescope to the latest observed entry, by Lemma~\ref{lem:telescope} and
$\hat C_{i,\,n-1-i} = C_{i,\,n-1-i}$ (the empty product).
\end{proof}

\begin{lemma}[Fitted column sums of cumulatives]\label{lem:colCum}\lean{VerifiedReserving.sum_CLcum_obsCol_aux, VerifiedReserving.sum_CLcum_obsCol, VerifiedReserving.sum_CLcum_contributors}\leanok
\uses{def:CLfit, def:S}
If every development factor $\hat f_j$, $j < n-1$, is nonzero, then for every $k \le n-1$
\[
  \sum_{i \le n-1-k} \hat C_{i,k} = \sum_{i \le n-1-k} C_{i,k},
\]
and, one year earlier over the same accident years,
$\sum_{i \le n-2-k} \hat C_{i,k} = S_k$.
\end{lemma}
\begin{proof}\leanok
Downward induction on $k$. At $k = n-1$ the only observed cell is the corner, where the
fit is the observation. For the step, the cells strictly above the diagonal satisfy
$\hat C_{i,k} = \hat C_{i,k+1}/\hat f_k$, so their sum is $T_k/\hat f_k = S_k$ by the
induction hypothesis and the definition $\hat f_k = T_k/S_k$; adding the diagonal cell,
where the fit is the observation, gives the claim. Nonvanishing of $\hat f_k$ forces
$T_k \ne 0$ and $S_k \ne 0$, which is the only place a hypothesis is used.
\end{proof}

\begin{theorem}[Column totals]\label{thm:odpCol}\lean{VerifiedReserving.chainLadder_fitted_column_totals}\leanok
\uses{lem:colCum, def:score}
If every development factor $\hat f_j$, $j < n-1$, is nonzero, then for every
development year $k$
\[
  \sum_{i \le n-1-k} \hat X_{i,k} = \sum_{i \le n-1-k} X_{i,k}:
\]
the chain-ladder fit reproduces every observed column total. This is the second
marginal-sum family of Mack (1991), Section 2, and the column score equation of Renshaw
and Verrall (1998), Section 3.
\end{theorem}
\begin{proof}\leanok
For $k = 0$ the incrementals are the cumulatives and this is Lemma~\ref{lem:colCum}. For
$k \ge 1$ subtract the two displays of Lemma~\ref{lem:colCum}: the first over the
observed part of column $k$, the second at $k-1$ over the same accident years.
\end{proof}

\begin{theorem}[Chain ladder solves the ODP score equations]\label{thm:odpScore}\lean{VerifiedReserving.chainLadder_scoreEquations}\leanok
\uses{thm:odpRow, thm:odpCol}
Under the same hypothesis, the chain-ladder fitted incrementals satisfy the whole
system of Definition~\ref{def:score}.
\end{theorem}
\begin{proof}\leanok Theorems~\ref{thm:odpRow} and~\ref{thm:odpCol}. \end{proof}

\section{Uniqueness: the marginal-sum solution is chain ladder}

Mack (1991), Section 2, states the converse: the marginal-sum estimates of the
multiplicative model are the chain-ladder estimates. Positivity of the parameters
replaces the informal assumption that the triangle is a real one.

\begin{lemma}[The pattern ratios are the development factors]\label{lem:patternFhat}\lean{VerifiedReserving.sum_mult_patternCum_aux, VerifiedReserving.sum_mult_patternCum, VerifiedReserving.patternCum_mul_fhat, VerifiedReserving.fhat_pos_of_mult}\leanok
\uses{def:mult, def:score, def:fhat}
Let $a_i, b_k > 0$ with $\sum_{k<n} b_k = 1$ satisfy the score equations. Then for every
$k \le n-1$
\[
  \Big(\sum_{i \le n-1-k} a_i\Big) B_k \text{ splits as } \sum_{i \le n-1-k} C_{i,k},
  \qquad\text{and}\qquad B_k \hat f_k = B_{k+1} \ (k < n-1),
\]
so every $\hat f_k$ is positive.
\end{lemma}
\begin{proof}\leanok
Downward induction on $k$, mirroring Lemma~\ref{lem:colCum}. The row equation for
accident year $i$ reads $a_i B_{n-1-i} = C_{i,\,n-1-i}$; at $i = 0$, with $B_{n-1} = 1$,
it starts the induction. In the step, write $A = \sum_{i \le n-2-k} a_i$. The induction
hypothesis gives $A B_{k+1} = T_k$ and the column equation at $k+1$ gives
$A b_{k+1} = T_k - S_k$, so $A B_k = S_k$; adding the diagonal cell, whose row equation
is $a_{n-1-k} B_k = C_{n-1-k,k}$, closes the induction. Dividing the two displays,
$\hat f_k = T_k/S_k = B_{k+1}/B_k$.
\end{proof}

\begin{lemma}[Fitted cumulatives agree]\label{lem:multCum}\lean{VerifiedReserving.patternCum_mul_prod, VerifiedReserving.mult_cum_eq_CLcum}\leanok
\uses{lem:patternFhat, def:CLfit}
Under the hypotheses of Lemma~\ref{lem:patternFhat}, $a_i B_k = \hat C_{i,k}$ on every
observed cell.
\end{lemma}
\begin{proof}\leanok
Telescoping $B_k \prod_{j=k}^{d-1}\hat f_j = B_d$ with $d = n-1-i$, and the row equation
$a_i B_d = C_{i,d}$; the product is nonzero because each factor is.
\end{proof}

\begin{theorem}[Mack's uniqueness result]\label{thm:odpUnique}\lean{VerifiedReserving.multFit_eq_CLincr}\leanok
\uses{lem:multCum}
A multiplicative model $m_{i,k} = a_i b_k$ with $a_i, b_k > 0$ and $\sum_{k<n} b_k = 1$
that satisfies the marginal-sum equations coincides with the chain-ladder fit cell by
cell on the observed triangle: $a_i b_k = \hat X_{i,k}$. With
Theorem~\ref{thm:odpScore} this is the classical equivalence: the over-dispersed
Poisson generalized linear model and the chain-ladder method give the same fitted
values, hence the same development pattern and the same reserves.
\end{theorem}
\begin{proof}\leanok
Differencing Lemma~\ref{lem:multCum} at $k$ and $k-1$, since
$B_k - B_{k-1} = b_k$.
\end{proof}
